\setcounter{theorem}{3}
\begin{theorem}\label{thm:main-sg-intro}
There exists a universal constant $c$ such that the following holds. Let $\widetilde{\mathcal Q}=\{Q_i\}_{i\in\{1,\ldots,m\}}\subset\mathbb C[x_1,\ldots,x_n]$ be a finite set of pairwise linearly independent homogeneous polynomials, such that every $Q_i\in\widetilde{\mathcal Q}$ is either irreducible or a square of a linear form. Assume that, for every $i\ne j$, whenever $Q_i$ and $Q_j$ vanish then so does $\prod_{k\in\{1,\ldots,m\}\setminus\{i,j\}}Q_k$. Then, $\dim(\operatorname{span}\{\mathcal Q\})\leq c$.
\end{theorem}
\begin{theorem}\label{thm:structure}\label{thm:structure-intro}
Let $\{Q_k\}_{k\in\mathcal K},A,B$ be homogeneous polynomials of degree2 such that $\prod_{k\in\mathcal K}Q_k\in\sqrt{\langle A,B\rangle}$. Then one of the following cases hold:
\begin{enumerate}[label={(\roman*)}]
\item\label{case:span} There is $k\in\mathcal K$ such that $Q_k$ is in the linear span of $A,B$.
\item\label{case:rk1} There exists a non trivial linear combination of the form $\alpha A+\beta B=c\cdot d$ where $c$ and $d$ are linear forms.
\item\label{case:2} There exist two linear forms $c$ and $d$ such that when setting $c=d=0$ we get that $A,B$ and one of $\{Q_k\}_{k\in\mathcal K}$ vanish.
\end{enumerate}
\end{theorem}
From now on, to ease notations, we use Theorem5i, Theorem5ii or Theorem5iii to describe different cases of Theorem5. The CCC formal version p8:5 refers the complete proof to its full version. The complete published author proof of precisely this structural statement is included below from journal Theorem3.1, pp11--16, CC BY4.0; the CCC main proof retains its own parameters and notation.
